% مشقوں کا مجموعہ — سبق ۲
% سبق کا تعین اس فائل کے نام سے ہوتا ہے۔

\begin{exo}[ٹھنڈی برف سے کھولتے پانی تک: مرتبۂ مقدار]{chauffage-glace-eau-ebullition}
\seoready{true}
\keywords{کیلوری میٹری, حالت کی تبدیلی, پگھلاؤ, تبخیر, پوشیدہ حرارت, مرتبۂ مقدار}

\textbf{عددی معلومات (فضائی دباؤ پر):}
\[
\begin{aligned}
c_{\mathrm{i}} &= 2{,}1\times10^3\,\text{J}\!\cdot\!\text{kg}^{-1}\!\cdot\!\text{K}^{-1},
& L_f &= 3{,}34\times10^5\,\text{J}\!\cdot\!\text{kg}^{-1},\\
c_{\mathrm{w}} &= 4{,}18\times10^3\,\text{J}\!\cdot\!\text{kg}^{-1}\!\cdot\!\text{K}^{-1},
& L_v &= 2{,}26\times10^6\,\text{J}\!\cdot\!\text{kg}^{-1}.
\end{aligned}
\]
فضائی دباؤ پر ایک کلوگرام برف کو، جس کا ابتدائی درجۂ حرارت $-100\,{}^\circ\text{C}$ ہے، بتدریج گرم کیا جاتا ہے۔

\begin{enumerate}
  \item برف کو $-100\,{}^\circ\text{C}$ سے $0\,{}^\circ\text{C}$ تک لانے کے لیے درکار توانائی نکالیں۔
  \item $0\,{}^\circ\text{C}$ پر برف کو مکمل طور پر پگھلانے کے لیے درکار توانائی نکالیں۔
  \item مائع پانی کو $0\,{}^\circ\text{C}$ سے $100\,{}^\circ\text{C}$ تک گرم کرنے کی توانائی نکالیں۔
  \item $100\,{}^\circ\text{C}$ پر پانی کو مکمل بخارات میں بدلنے کے لیے درکار اضافی توانائی نکالیں۔
  \item کس مرحلے میں سب سے زیادہ توانائی درکار ہے؟ ایک کلوگرام پانی کو $0\,{}^\circ\text{C}$ سے $100\,{}^\circ\text{C}$ تک گرم کرنے کی توانائی کا موازنہ دس میٹر کی بلندی سے گرتی ہوئی کمیت $m$ کی ثقلی مخزونی توانائی سے کریں۔ اتنی ہی توانائی کے لیے کتنی کمیت $m$ گرانی ہوگی؟
  \item اب الٹا عمل دیکھیے۔ آبی بخارات کی کمیت $m_v$ گاڑی کی سامنے والی شیشے پر تکثیف پاتی ہے اور حاصل شدہ مائع پانی بعد میں ٹھنڈا نہیں ہوتا۔ تکثیف کے دوران خارج شدہ حرارت $Q_{\mathrm{rel}}$ کا اظہار لکھیں، پھر $m_v=1{,}0\times10^{-2}\,\text{kg}$ کے لیے حساب کریں۔ شیشے کے درجۂ حرارت پر $L_v=2{,}45\times10^6\,\text{J}\!\cdot\!\text{kg}^{-1}$ لیں۔
\end{enumerate}

\begin{indication}
حالت بدلے بغیر گرم کرنے کے لیے $Q=mc\Delta T$ اور مستقل درجۂ حرارت پر حالت کی تبدیلی کے لیے $Q=mL$ استعمال کریں۔ بلندی $h$ پر کمیت $m$ کی ثقلی مخزونی توانائی $E_p=mgh$ ہے؛ $g=9{,}81\,\text{m}\!\cdot\!\text{s}^{-2}$ لیں۔
\end{indication}

\begin{solution}
\textbf{1.} برف کو نقطۂ پگھلاؤ تک گرم کرنے کے لیے
\[
Q_1=mc_{\mathrm{i}}\Delta T
=1{,}0\times2{,}1\times10^3\times100
=2{,}1\times10^5\,\text{J}.
\]

\textbf{2.} پگھلاؤ کے دوران درجۂ حرارت $0\,{}^\circ\text{C}$ رہتا ہے:
\[
Q_2=mL_f=1{,}0\times3{,}34\times10^5=3{,}34\times10^5\,\text{J}.
\]

\textbf{3.} مائع پانی کو $0$ سے $100\,{}^\circ\text{C}$ تک گرم کرنے کے لیے
\[
Q_3=mc_{\mathrm{w}}\Delta T
=1{,}0\times4{,}18\times10^3\times100
=4{,}18\times10^5\,\text{J}.
\]

\textbf{4.} مکمل تبخیر کے لیے مزید
\[
Q_4=mL_v=1{,}0\times2{,}26\times10^6=2{,}26\times10^6\,\text{J}
\]
درکار ہیں۔ صرف یہ مرحلہ ہی چند میگاجول کے مرتبے کی توانائی لیتا ہے۔

\textbf{5.} تبخیر واضح طور پر سب سے زیادہ توانائی طلب مرحلہ ہے:
\[
Q_4=2{,}26\times10^6\,\text{J}
>Q_3=4{,}18\times10^5\,\text{J}
>Q_2=3{,}34\times10^5\,\text{J}
>Q_1=2{,}1\times10^5\,\text{J}.
\]
تبخیر کو پانی $0$ سے $100\,{}^\circ\text{C}$ تک گرم کرنے کے مقابلے میں تقریباً $(2{,}26\times10^6)/(4{,}18\times10^5)\approx5{,}4$ گنا توانائی درکار ہے۔

$h=10\,\text{m}$ سے گرتی کمیت کے لیے $E_p=mgh$۔ شرط $mgh=Q_3$ سے
\[
m=\frac{Q_3}{gh}
=\frac{4{,}18\times10^5}{9{,}81\times10}
\approx4{,}26\times10^3\,\text{kg}.
\]
یعنی ایک کلوگرام پانی کو $0$ سے $100\,{}^\circ\text{C}$ تک گرم کرنے جتنی توانائی کے لیے تقریباً $4{,}3$ ٹن کو دس میٹر سے گرانا ہوگا! اسی لیے گھریلو توانائی کے استعمال میں پانی گرم کرنا اہم حصہ ہے۔ پانی کی مخصوص حرارتی گنجائش عام دھاتوں سے بھی بہت زیادہ ہے؛ مثلاً تانبے کی تقریباً دس گنا (اگلی مشقیں دیکھیے)۔

\textbf{6.} تکثیف تبخیر کا الٹا عمل ہے؛ بخارات شیشے کو پوشیدہ حرارت دیتے ہیں:
\[
Q_{\mathrm{rel}}=m_vL_v.
\]
$m_v=1{,}0\times10^{-2}\,\text{kg}$ کے لیے
\[
Q_{\mathrm{rel}}=1{,}0\times10^{-2}\times2{,}45\times10^6
=2{,}45\times10^4\,\text{J}.
\]
یوں تھوڑی سی بھاپ کی تکثیف بھی قابلِ لحاظ حرارت خارج کر سکتی ہے، جسے شیشہ اور اس کا ماحول جذب کرتے ہیں۔
\end{solution}
\end{exo}

\begin{exo}[ایک بڑے درخت کا تبخیر و تعرق: قدرتی ایئر کنڈیشنر؟]{odg-evapotranspiration-arbre}
\seoready{true}
\keywords{مرتبۂ مقدار, تبخیر و تعرق, درخت, پوشیدہ حرارت, تبریدی طاقت, ایئر کنڈیشنر, COP}

گرم اور خشک دن میں وافر پانی پانے والا ایک بڑا درخت تقریباً $500\,\text{L}=5{,}0\times10^{-1}\,\text{m}^3$ پانی تبخیر و تعرق سے خارج کر سکتا ہے۔ چونکہ تعرق زیادہ تر روشنی میں ہوتا ہے، اسے بارہ گھنٹوں میں پھیلا ہوا فرض کریں۔ درخت کا پھیلاؤ زمین پر $A_{\mathrm{tree}}=50\,\text{m}^2$ رقبہ ڈھانپتا ہے۔ فضائی دباؤ پر پانی کی کثافت اور مخصوص پوشیدہ حرارتِ تبخیر یہ ہیں:
\[
\rho_{\mathrm{w}}=1{,}0\times10^3\,\text{kg}\!\cdot\!\text{m}^{-3},
\qquad L_v=2{,}45\times10^6\,\text{J}\!\cdot\!\text{kg}^{-1}.
\]

\begin{enumerate}
  \item دن میں خارج ہونے والے پانی کی کمیت اور اسے بخارات بنانے کی توانائی نکالیں۔
  \item بتائیں کہ اس سے ٹھنڈک کیوں پیدا ہوتی ہے۔
  \item تعرق کے بارہ فعال گھنٹوں میں درخت کی فی مربع میٹر اوسط تبریدی طاقت نکالیں۔
  \item موازنے کے لیے $P_{\mathrm{elec}}=2{,}0\times10^3\,\text{W}$ برقی طاقت اور ضریبِ کارکردگی $\mathrm{COP}=3{,}0$ والا ایئر کنڈیشنر $A_{\mathrm{room}}=20\,\text{m}^2$ کے کمرے کو ٹھنڈا کرتا ہے۔ تعریفاً
\[
\mathrm{COP}=\frac{P_{\mathrm{cool}}}{P_{\mathrm{elec}}},
\]
جہاں $P_{\mathrm{cool}}$ کمرے سے نکالی گئی حرارتی طاقت ہے۔ $P_{\mathrm{cool}}$ اور فی مربع میٹر تبریدی طاقت نکال کر درخت سے موازنہ کریں۔
  \item فی اکائی رقبہ طاقتوں کا مرتبۂ مقدار یکساں ہونے کے باوجود یہ نتیجہ کیوں نہیں نکلتا کہ درخت اور ایئر کنڈیشنر بالکل ایک جیسی محسوس شدہ ٹھنڈک دیتے ہیں؟
  \item کیا بہت سے اندرونی پودوں سے گھر کو ٹھنڈا کیا جا سکتا ہے؟ (عمومی معلومات: جواب حیاتیاتی عوامل پر منحصر ہے۔)
\end{enumerate}

\begin{indication}
$m=\rho V$، $Q_{\mathrm{evap}}=mL_v$ اور $P=Q/\tau$ استعمال کریں۔ ایئر کنڈیشنر کے لیے $P_{\mathrm{cool}}=\mathrm{COP}\,P_{\mathrm{elec}}$۔
\end{indication}

\begin{solution}
\textbf{1.} $500\,\text{L}=5{,}0\times10^{-1}\,\text{m}^3$، لہٰذا
\[
m=\rho_{\mathrm{w}}V=1{,}0\times10^3\times5{,}0\times10^{-1}=5{,}0\times10^2\,\text{kg},
\]
اور
\[
Q_{\mathrm{evap}}=mL_v=5{,}0\times10^2\times2{,}45\times10^6=1{,}225\times10^9\,\text{J}.
\]
یہ تقریباً ایک گیگاجول ہے۔

\textbf{2.} تبخیر حرارت گیر حالت کی تبدیلی ہے۔ مطلوبہ $Q_{\mathrm{evap}}$ پتوں، زمین اور آس پاس کی ہوا سے آتی ہے؛ حرارت کے اس اخراج سے ان کا درجۂ حرارت کم ہوتا ہے۔

\textbf{3.} بارہ گھنٹے $\tau=12\times3600=4{,}32\times10^4\,\text{s}$ ہیں، اس لیے
\[
P_{\mathrm{tree}}=\frac{Q_{\mathrm{evap}}}{\tau}
=\frac{1{,}225\times10^9}{4{,}32\times10^4}
\approx2{,}84\times10^4\,\text{W},
\]
اور $50\,\text{m}^2$ پر
\[
\frac{P_{\mathrm{tree}}}{A_{\mathrm{tree}}}\approx5{,}7\times10^2\,\text{W}\!\cdot\!\text{m}^{-2}.
\]

\textbf{4.}
\[
P_{\mathrm{cool}}=\mathrm{COP}\,P_{\mathrm{elec}}=3{,}0\times2{,}0\times10^3=6{,}0\times10^3\,\text{W},
\]
اور
\[
\frac{P_{\mathrm{cool}}}{A_{\mathrm{room}}}=\frac{6{,}0\times10^3}{20}=3{,}0\times10^2\,\text{W}\!\cdot\!\text{m}^{-2}.
\]
اس نمونے میں درخت کا تبخیر و تعرق عام ایئر کنڈیشنر سے تقریباً دو گنا طاقتور ہے۔

\textbf{5.} موازنہ صرف توانائی کے بہاؤ کا ہے۔ ایئر کنڈیشنر بند اور قابو شدہ حجم ٹھنڈا کرتا ہے، جب کہ درخت کی بھاپ فضا میں پھیلتی اور ہوا گرم ہوا واپس لاتی ہے۔ محسوس اثر ہواداری، نمی، محیطی درجۂ حرارت اور درخت سے فاصلے پر منحصر ہے۔ تعرق بھی دھوپ، ہوا، نمی، نوع اور دستیاب پانی کے ساتھ بدلتا اور خشک سالی میں بہت کم ہو سکتا ہے۔ دوسری طرف درخت سایہ بھی دیتا ہے۔ اس لیے یہ حساب مرتبۂ مقدار کا موازنہ ہے، حرارتی آسائش کی عین مساوات نہیں۔

\textbf{6.} عملی طور پر اندرونی پودوں کی ٹھنڈک بہت معمولی ہوگی۔ بیشتر پودوں میں روشنی روزن کھولتی ہے جن سے بھاپ نکلتی ہے؛ کم اندرونی روشنی میں یہ کم کھلتے ہیں۔ کم ہوادار کمرے میں بڑھتی نمی تبخیر کو بھی سست کرتی ہے۔ بہت سے پودے ہلکی مقامی تبخیری ٹھنڈک تو دے سکتے ہیں، مگر ایئر کنڈیشنر کا بدل نہیں۔ نم ہوا کو باہر نکالنے کے لیے ہوا بدلنا لازم ہوگا۔ تیز مصنوعی روشنی تعرق بڑھا سکتی ہے لیکن اس کی برقی توانائی بالآخر تقریباً پوری حرارت بن کر کمرے ہی میں آتی ہے۔ کیکٹس جیسے خشک ماحول کے CAM پودے استثنا ہیں: ان کے روزن عموماً رات کو کھلتے ہیں۔
\end{solution}
\end{exo}

\begin{exo}[پانی کی دو کمیتوں کا اختلاط]{calorimetrie-melange-eau}
\seoready{true}
\keywords{کیلوری میٹری, اختلاط, حرارتی گنجائش, مخصوص حرارت, جوزف بلیک, توازنی درجۂ حرارت}

$T_1=80\,{}^\circ\text{C}$ پر پانی کی کمیت $m_1=0{,}200\,\text{kg}$ ایک ایسے ظرف میں ڈالی جاتی ہے جس میں پہلے سے $T_2=15\,{}^\circ\text{C}$ پر $m_2=0{,}300\,\text{kg}$ پانی ہے۔ ظرف مثالی حرارت پیما ہے: باہر حرارت کا نقصان اور ظرف کی اپنی حرارتی گنجائش نظرانداز ہیں۔ یاد رہے
\[
Q=mc\Delta T,
\]
جہاں مائع پانی کی مخصوص حرارتی گنجائش $c=4{,}18\times10^3\,\text{J}\!\cdot\!\text{kg}^{-1}\!\cdot\!\text{K}^{-1}$ ہے۔

\begin{enumerate}
  \item واضح کریں کہ گرم پانی کی دی ہوئی پوری حرارت ٹھنڈا پانی لیتا ہے؛ $m_1$، $m_2$، $c$، $T_1$، $T_2$ اور توازنی درجۂ حرارت $T_{eq}$ کی مساوات لکھیں۔
  \item $T_{eq}$ کا جبری اظہار اور عددی قدر نکالیں۔
  \item اختلاط میں گرم پانی کی دی ہوئی حرارت $Q$ نکالیں۔
  \item کیا ایک ہی مادے کے ایسے اختلاط سے $c$ کی قدر معلوم ہو سکتی ہے؟ وجہ بتائیں۔
\end{enumerate}

\begin{indication}
حرارت پیما بند اور سخت ہے، اس لیے کل داخلی توانائی $U_1+U_2$ محفوظ ہے: ایک کی دی ہوئی حرارت دوسرے کو الٹی علامت کے ساتھ ملتی ہے۔
\end{indication}

\begin{solution}
\textbf{1.} بند حرارت پیما میں گرم اور ٹھنڈے پانی کے نظام کی کل داخلی توانائی محفوظ ہے۔ ہر کمیت پر ابتدائی حالت سے توازن تک $Q=mc\Delta T$ لگانے سے
\[
m_1c(T_{eq}-T_1)+m_2c(T_{eq}-T_2)=0.
\]

\textbf{2.} ایک ہی مادہ ہونے سے $c$ منسوخ ہو جاتا ہے:
\[
T_{eq}=\frac{m_1T_1+m_2T_2}{m_1+m_2}
=\frac{0{,}200\times80+0{,}300\times15}{0{,}200+0{,}300}
=41\,{}^\circ\text{C}.
\]

\textbf{3.}
\[
Q=m_1c(T_1-T_{eq})
=0{,}200\times4{,}18\times10^3\times(80-41)
\approx3{,}26\times10^4\,\text{J}.
\]
ٹھنڈے پانی کو عین اتنی ہی حرارت ملتی ہے:
\[
m_2c(T_{eq}-T_2)=0{,}300\times4{,}18\times10^3\times26\approx3{,}26\times10^4\,\text{J}.
\]

\textbf{4.} نہیں۔ ایک ہی مادے کی دونوں کمیتوں کے توانائی موازنے میں یکساں $c$ دونوں حدوں کا عامل ہے اور منسوخ ہو جاتا ہے۔ توازنی درجۂ حرارت $c$ پر نہیں بلکہ صرف کمیتوں سے وزن شدہ ابتدائی درجۂ حرارتوں کے اوسط پر منحصر ہے۔ طریقۂ اختلاط سے مخصوص حرارتی گنجائش معلوم کرنے کے لیے اگلی مشق کی طرح معلوم حرارتی گنجائش والا مادہ شامل کرنا ہوگا۔
\end{solution}
\end{exo}

\begin{exo}[طریقۂ اختلاط سے دھات کی مخصوص حرارتی گنجائش کا تعین]{calorimetrie-methode-melanges}
\seoready{true}
\keywords{کیلوری میٹری, طریقۂ اختلاط, مخصوص حرارتی گنجائش, حرارت پیما, مخصوص حرارت, آبی مساوی}

جیسے جوزف بلیک اٹھارہویں صدی میں کرتے تھے، ہم \emph{طریقۂ اختلاط} سے نامعلوم دھات کی مخصوص حرارتی گنجائش $c_m$ ناپنا چاہتے ہیں۔ $m=0{,}150\,\text{kg}$ نمونے کو $T_m=95\,{}^\circ\text{C}$ تک گرم کر کے جلدی سے اس حرارت پیما میں ڈالتے ہیں جس میں $T_e=18\,{}^\circ\text{C}$ پر $m_e=0{,}250\,\text{kg}$ پانی ہے؛ $c_e=4{,}18\times10^3\,\text{J}\!\cdot\!\text{kg}^{-1}\!\cdot\!\text{K}^{-1}$۔ حرارتی توازن پر $T_{eq}=22\,{}^\circ\text{C}$ ملتا ہے۔ پہلے ظرف (دیواریں، ہلانے والا اور تھرمامیٹر) کی اپنی حرارتی گنجائش نظرانداز کریں۔

\begin{enumerate}
  \item پچھلی مشق کی طرح دھات اور پانی کے بند نظام میں توانائی کی بقا لکھیں، جہاں $c_m$ واحد نامعلوم ہو۔
  \item $c_m$ کا جبری اظہار اور عددی قدر نکالیں۔ یہ تقریباً کس عام دھات سے مطابقت رکھتی ہے؟ حوالہ: المونیم $\approx9{,}0\times10^2\,\text{J}\!\cdot\!\text{kg}^{-1}\!\cdot\!\text{K}^{-1}$، لوہا $\approx4{,}5\times10^2$، تانبا $\approx3{,}9\times10^2$، سیسہ $\approx1{,}3\times10^2$۔
  \item حقیقت میں ظرف بھی دھات کی دی ہوئی حرارت کا کچھ حصہ لیتا ہے۔ اسے \emph{آبی مساوی} $\mu=1{,}5\times10^{-2}\,\text{kg}$ سے نمونہ بنائیں: حرارت پیما گویا پانی کی اضافی کمیت $\mu$ ہے۔ $\mu$ شامل کر کے $c_m$ دوبارہ نکالیں اور اصلاح کی سمت پر تبصرہ کریں (کیا درست قدر سوال ۲ سے بڑی ہے یا چھوٹی؟)۔
  \item تجربے میں نمونے کو \emph{جلدی} حرارت پیما میں ڈالنا اور پیمائش کے دوران محیطی ہوا سے اچھی طرح حرارتی طور پر الگ رکھنا کیوں ضروری ہے؟
\end{enumerate}

\begin{indication}
دھات $Q_m=mc_m(T_m-T_{eq})$ حرارت دیتی ہے، جسے پانی (اور سوال ۳ میں اضافی آبی کمیت $\mu$ کے مساوی حرارت پیما) مکمل طور پر لیتا ہے: $Q_m=(m_e+\mu)c_e(T_{eq}-T_e)$۔
\end{indication}

\begin{solution}
\textbf{1.} بند نظام میں دھات $T_m$ سے $T_{eq}$ تک ٹھنڈی ہو کر جو حرارت دیتی ہے، پانی $T_e$ سے $T_{eq}$ تک گرم ہو کر پوری لیتا ہے:
\[
mc_m(T_m-T_{eq})=m_ec_e(T_{eq}-T_e).
\]

\textbf{2.}
\[
c_m=\frac{m_ec_e(T_{eq}-T_e)}{m(T_m-T_{eq})}
=\frac{0{,}250\times4{,}18\times10^3\times(22-18)}{0{,}150\times(95-22)}
\approx3{,}82\times10^2\,\text{J}\!\cdot\!\text{kg}^{-1}\!\cdot\!\text{K}^{-1}.
\]
یہ تانبے کی قدر ($\approx3{,}9\times10^2\,\text{J}\!\cdot\!\text{kg}^{-1}\!\cdot\!\text{K}^{-1}$) کے بہت قریب ہے۔

\textbf{3.} اب پانی اور ظرف دونوں حرارت لیتے ہیں:
\[
c_m=\frac{(m_e+\mu)c_e(T_{eq}-T_e)}{m(T_m-T_{eq})}
=\frac{(0{,}250+0{,}015)\times4{,}18\times10^3\times4}{0{,}150\times73}
\approx4{,}05\times10^2\,\text{J}\!\cdot\!\text{kg}^{-1}\!\cdot\!\text{K}^{-1}.
\]
اصلاح $c_m$ کو قدرے بڑھاتی ہے۔ سوال ۲ میں آبی مساوی نظرانداز کرنے سے دھات کی اصل دی ہوئی حرارت—جس کا کچھ حصہ ظرف گرم کرنے میں گیا—اور نتیجتاً $c_m$ کم سمجھا گیا تھا۔

\textbf{4.} طریقہ اس مفروضے پر قائم ہے کہ پیمائش کے دوران نظام بند ہے۔ نمونہ جلدی ڈالنے سے پانی تک پہنچنے سے پہلے ہوا سے حرارت کے تبادلے کا وقت کم ہوتا ہے، اور اچھی حرارتی علیحدگی توازن قائم ہوتے وقت باہر نقصان گھٹاتی ہے۔ کوئی غیر محسوب حرارتی تبادلہ سوال ۱ کا موازنہ اور اخذ شدہ $c_m$ غلط کر دے گا۔ بلیک اور بعد میں لاووازیے و لاپلاس نے اپنے برفی حرارت پیما میں اسی نوع کی تجرباتی احتیاط کی تھی۔
\end{solution}
\end{exo}

\begin{exo}[غذائی کیلوریاں اور میٹابولک طاقت]{calorimetrie-calories-alimentaires}
\seoready{true}
\keywords{کیلوری, کلوکیلوری, غذائی کیلوری, میکانیکی مساوی, میٹابولک طاقت, اکائیاں}

غذائی لیبل کے مطابق ایک بالغ روزانہ اوسطاً تقریباً $2000\,\text{kcal}$ لیتا ہے۔ کلوکیلوری ایک غیر SI اکائی ہے جو اب بھی غذائیت میں استعمال ہوتی ہے؛ اس کی تبدیلی $1\,\text{kcal}=4{,}18\times10^3\,\text{J}$ ہے۔

\begin{enumerate}
  \item اس یومیہ توانائی کو جول میں بدلیں۔
  \item فرض کریں کہ توانائی ۲۴ گھنٹوں میں یکساں استعمال ہوتی ہے؛ بالغ کی اوسط میٹابولک طاقت واٹ میں نکالیں۔
  \item ایک توانائی بار پر ``$210\,\text{kcal}$'' لکھا ہے۔ اگر یہی توانائی مکمل طور پر حرارت میں بدل جائے تو $20\,{}^\circ\text{C}$ کے کتنی کمیت والے پانی کو نقطۂ ابال ($100\,{}^\circ\text{C}$) تک لایا جا سکتا ہے؟ $c_{\mathrm{w}}=4{,}18\times10^3\,\text{J}\!\cdot\!\text{kg}^{-1}\!\cdot\!\text{K}^{-1}$ لیں۔
  \item آپ کے خیال میں انسانی جسم اس توانائی کو حقیقت میں کن صورتوں میں استعمال کرتا ہے؟ (کیفی جواب)
\end{enumerate}

\begin{indication}
سوال ۲ کے لیے $\mathcal P=E/\Delta t$ استعمال کریں۔ سوال ۳ میں پہلے لیبل کی توانائی جول میں بدلیں، پھر $Q=mc\Delta T$ سے $m$ الگ کریں۔
\end{indication}

\begin{solution}
\textbf{1.} SI اکائیوں میں
\[
E=2000\times4{,}18\times10^3=8{,}36\times10^6\,\text{J}.
\]

\textbf{2.}
\[
\mathcal P=\frac{8{,}36\times10^6}{8{,}64\times10^4}\approx97\,\text{W}.
\]
پورے دن میں انسان اوسطاً تقریباً ایک روایتی روشن بلب جتنی طاقت استعمال کرتا ہے۔ یہ حیران کن مرتبۂ مقدار غذائی کیلوری کو عام طبیعی اکائیوں میں بدلنے کی افادیت دکھاتا ہے۔

\textbf{3.} بار کی توانائی
\[
Q=210\times4{,}18\times10^3\approx8{,}78\times10^5\,\text{J}.
\]
$\Delta T=80\,\text{K}$ کے ساتھ
\[
m=\frac{Q}{c\Delta T}
=\frac{8{,}78\times10^5}{4{,}18\times10^3\times80}
\approx2{,}63\,\text{kg}.
\]
یعنی مرتبۂ مقدار میں ایک توانائی بار تقریباً $2{,}5\,\text{kg}$ نلکے کے پانی کو ابال تک پہنچانے کے برابر ہے۔

\textbf{4.} جسم توانائی کو غائب کر کے ``خرچ'' نہیں کرتا بلکہ غذا کی کیمیائی توانائی کو بدلتا ہے۔ کچھ حصہ عضلات کے میکانیکی کام، اعضا کے افعال، خلیاتی برقی میلان برقرار رکھنے، سالمات کی نقل و حمل اور بافتوں کی تجدید کے کیمیائی تالیف میں لگتا ہے۔ کچھ گلائیکوجن اور چربی میں کیمیائی طور پر ذخیرہ ہو سکتا ہے۔ بالآخر استعمال شدہ توانائی کا بیشتر حصہ حرارت بن کر منتشر ہوتا ہے، جو ماحول کو منتقل ہونے سے پہلے جسمانی درجۂ حرارت برقرار رکھنے میں مدد دیتا ہے۔ تھوڑا حصہ فضلے میں باقی کیمیائی توانائی کے طور پر نکلتا ہے۔
\end{solution}
\end{exo}
